\documentclass[10pt,a4paper]{amsart}
\usepackage[margin=19mm]{geometry}
\usepackage{amsmath,amssymb,mathtools}
\usepackage[scr=rsfs,cal=euler]{mathalfa}
\usepackage{xfrac}
\usepackage[unicode,hidelinks,bookmarks=false]{hyperref}
\newcommand{\ie}{i.e.\ }
\newcommand{\cf}{cf.\ }
\newcommand{\resp}{respectively\ }
\newcommand{\Subsection}{\subsection}
\providecommand{\nobreakdash}{\nobreak\hbox{-}}
\setlength{\emergencystretch}{2em}
\multlinegap0pt
\usepackage{bm}
\usepackage{bbm}
\usepackage{dsfont}
\usepackage{xspace}
\usepackage{cancel}
\usepackage{mathtools}
\usepackage{amsthm}
\makeatletter
\@ifundefined{assumption}{\newtheorem{assumption}{Assumption}}{}
\@ifundefined{lemma}{\newtheorem{lemma}{Lemma}}{}
\@ifundefined{remark}{\newtheorem{remark}{Remark}}{}
\@ifundefined{theorem}{\newtheorem{theorem}{Theorem}}{}
\makeatother
\newcommand{\one}{{\bf 1}}
\title[Outlier eigenvalues and eigenvectors of generalized Wigner m]{Outlier eigenvalues and eigenvectors of generalized Wigner matrices with finite-rank perturbations}
\author{Bishakh Bhattacharya, Arijit Chakrabarty, Rajat Subhra Hazra}
\date{Source: arXiv:2601.10204v1 (CC BY 4.0)}
\begin{document}
\maketitle

\noindent\emph{Source and licence.} Outlier eigenvalues and eigenvectors of generalized Wigner matrices with finite-rank perturbations. Version arXiv:2601.10204v1 (CC BY 4.0), distributed under the Creative Commons Attribution 4.0 International licence, as recorded in the arXiv licence metadata for this exact version and shown on the abstract page https://arxiv.org/abs/2601.10204v1. Full rights notice: licenses/arxiv-2601.10204-CC-BY-4.0.txt.\par\smallskip

\noindent\textit{Editorial scope.} Counted unit: the Lemma whose source label is \texttt{EVec.pivot1} in the article (source0.tex lines 896--916, source label \texttt{EVec.pivot1}), delivered together with the article's own complete original proof of it (source0.tex lines 918--1074, one \texttt{proof} environment of 157 lines). No mathematics is written, repaired or re-derived: statement and proof are the article's own text line for line. Required context retained uncounted: eq:WignerTypeMatrix (lines 269--270); eq:momentcond (lines 276--277); assump:A1 (lines 292--298); assump:A2 (lines 303--305); assump:A3 (lines 311--317); lemma:NormBound (lines 533--536); Th.probedge (lines 559--563); Bound.W/mu (lines 594--600); eq:mu-mu0 (lines 691--693); EigenVector.1 (lines 747--752); EigenVector.3 (lines 828--843); EigenVector.3.intermd (lines 861--867); Remark:djlm (lines 887--894); EigenVector.3.exten (lines 889--893). Every definition, display and label of those lines is retained unchanged. Omitted from this package and not redistributed: 1--268, 271--275, 278--291, 299--302, 306--310, 318--532, 537--558, 564--593, 601--690, 694--746, 753--827, 844--860, 868--886, 894--895, 917--917, 1075--3140. Nothing inside a retained excerpt is omitted or altered: every retained body is delivered exactly as the article's lines are written, so no omission receipt of any kind accompanies this package. Citations: the retained excerpts carry no citation command at all, so no bibliography entry of the article is transcribed and no reference is redistributed. Reference adaptation: none was needed and none was made; every reference inside the delivered unit resolves inside it, so no unresolved reference or citation is produced. Printed slips kept literal: no printed word, symbol, hypothesis or number of the counted statement or of its proof was altered. Layout and class: the article's own class is \texttt{amsart} with packages including charter, mathrsfs, graphicx, geometry, amsmath, amsthm, amssymb, enumerate, fancyhdr, bbm, natbib, hyperref, mathtools, xcolor; the wrapper is the lane's pinned \texttt{amsart} editorial wrapper at 10pt on A4, which restates the theorem-like environments the retained text needs and adds the article's own macro definitions that the retained text needs (\texttt{\textbackslash one}), with extra wrapper packages (bm, bbm, dsfont, xspace, cancel, mathtools). The delivered numbering is the unit's own. Assets: no retained span contains a figure environment, a graphics-inclusion command, a PDF-page inclusion, a table or a TikZ picture; no image file is copied into this package, the package declares no asset dependency and no image file is redistributed with it. Licence: version arXiv:2601.10204v1 of this article is distributed under the Creative Commons Attribution 4.0 International licence, as recorded in the arXiv licence metadata for this exact version and shown on the abstract page https://arxiv.org/abs/2601.10204v1. A full rights notice is delivered at licenses/arxiv-2601.10204-CC-BY-4.0.txt. Characters: every retained excerpt is pure ASCII, so the delivered engine, XeLaTeX with its default Latin Modern text font, reports no missing character.
\begin{equation}\label{eq:WignerTypeMatrix} W_N(j,\ell) \;=\; X_{\,j \wedge \ell,\, j \vee \ell}\,,
\end{equation} where $\{X_{j,\ell}: 1 \le j \le \ell \le N\}$ is a collection of independent (not necessarily

that for all $m \ge 1$, \begin{equation}\label{eq:momentcond} \mathbb{E}[\,|X_{j,\ell}|^m\,] \;\le\; m^{\,B
m}\,. \end{equation} This condition (sometimes called a sub-exponential tail bound) is used to ensure

\begin{assumption}[{\bf Growth of \(\theta_N\)}]\label{assump:A1}
There exists $\xi > 4(B+1) \vee 8$ such that
\[
\sqrt{N}(\log N)^{\xi} \preceq \theta_N \preceq N
\]
for all large $N$.
\end{assumption}

\begin{assumption}[{\bf Orthogonality of \(h_i\)}]\label{assump:A2}
    The collection ${h_1,\ldots,h_k}$ is an orthonormal set of bounded functions in $L^2[0,1]$. 
\end{assumption}

\begin{assumption}[{\bf Regularity of $h_i$}]\label{assump:A3} Each $h_i: [0,1]\to\mathbb R$ is Lipschitz continuous: there exists $L_i < \infty$ such that

\[
        |h_i(x) - h_i(y)| \leq L_i |x - y|, \quad \forall x, y \in [0,1].
    \]

\end{assumption}

\begin{lemma}({\bf Norm bound for $W$})\label{lemma:NormBound} Let $W$ be the Wigner-type matrix defined in \eqref{eq:WignerTypeMatrix} satisfying \eqref{eq:momentcond} and with variance profile $f(x,y)$. Then \[
\|W\| = O_{hp}(\sqrt{N})\,.
\]
\end{lemma}

\begin{theorem}\label{Th.probedge}\textbf{First-Order Behavior of Eigenvalues.} Under Assumptions \eqref{assump:A1} and \eqref{assump:A2}, we have for every $1 \leq i \leq k$,
\[
\lambda_i(A) = \theta \alpha_i \left(1 + o_{\mathrm{hp}}(1)\right)\,.
\]
\end{theorem}

\begin{remark}\label{Bound.W/mu}
 From Theorem \ref{Th.probedge} we have $\mu = \lambda_i(A) \sim \theta \alpha_i$ for large $N$, with high probability. Meanwhile, Lemma \ref{lemma:NormBound} says $\|W\| = O_{hp}(\sqrt{N}) = o_{hp}(\theta)$ by Assumption \ref{assump:A1}. Combining these, we get
\begin{align*}
\left\|\frac{W}{\mu}\right\|=O_{hp}\left(\frac{\sqrt{N}}{\theta}\right) = o_{hp}(1)
\end{align*}
In particular, for large $N$, the operator $I - \frac{W}{\mu}$ is invertible with high probability, which justifies the expansion of the resolvent in a Neumann series.
\end{remark}

\begin{equation}\label{eq:mu-mu0}
    \mu-\mu_0=O_{hp}\left(\frac{\|Y_1\|}{\theta}+1\right).
    \end{equation}

\begin{lemma}\label{EigenVector.1}
With Assumptions \ref{assump:A1}-\ref{assump:A3} in place, as $N \to \infty$ we have %for $j \ne l$
\[
e_j^\prime\left(I - \frac{W}{\mu}\right)^{-m}e_l = \one(j =l) + o_{hp}(\frac{\sqrt{N}}{\theta}), \text{ for } m\in \{1,2\}.
\]
\end{lemma}

\begin{lemma}\label{EigenVector.3}
Under Assumptions \ref{assump:A1}-\ref{assump:A3}, as $N \to \infty$, for $j,l,p \ne i$, we have
\begin{align*}
\left[e_j^\prime\left(I - \frac{W}{\mu}\right)^{-1}e_i\right]
\left[e_l^\prime\left(I - \frac{W}{\mu}\right)^{-1}e_i\right]
\left[e_p^\prime\left(I - \frac{W}{\mu}\right)^{-2}e_p\right] 
= D_{lpp}S_2^{ji} + o_{hp}\left( \frac{\sqrt{N}}{\theta^2}\right),
\end{align*}
where
\begin{align*}
D_{lpp} &= \frac{3\,\mathbb E(e_j^\prime W^2 e_i) \mathbb E(e_l^\prime W^2 e_i)}{\mu_0^4}
+ \frac{\mathbb E(e_l^\prime W^2 e_i)}{\mu_0^2}\tilde S_3^{pp}
+ \frac{3\,\mathbb E(e_p^\prime W^2 e_p)}{\mu_0^2}S_3^{li} \\
&\quad + e_l^\prime e_i e_p^\prime e_p + e_p^\prime e_p S_2^{li} + S_3^{li}\tilde S_3^{pp}.
\end{align*}
\end{lemma}

\begin{align}
&\left[e_l^\prime\left(I - \frac{W}{\mu}\right)^{-1}e_i\right]\left[e_p^\prime\left(I - \frac{W}{\mu}\right)^{-2}e_p\right] \nonumber\\
&= D_{lpp} + \frac{e_l^\prime W e_i}{\mu}(e_p^\prime e_p)
+ o_{hp}\left(\left(1 + \frac{\|Y_1\|}{\theta} + \frac{\|Y_1\|^2}{\theta^2}\right)\frac{N}{\theta^3} + \frac{\sqrt{N}}{\theta^2}\right) \nonumber\\
&\quad + O_{hp}\left(\frac{N(1+\frac{\|Y_1\|}{\theta})}{\theta^3}\right)
+ O_{hp}\left(\frac{\sqrt{N}(\log N)^{\xi/2}}{\theta^2}\right),\label{EigenVector.3.intermd}
\end{align}

\begin{remark}\label{Remark:djlm}
It is clear that for $j,l,m \ne i$, our estimate in \eqref{EigenVector.3.intermd} simplifies as  
\begin{align}
 & \left[e_j^\prime\left(I - \frac{W}{\mu}\right)^{-1}e_i\right]\left[e_l^\prime\left(I - \frac{W}{\mu}\right)^{-2}e_m\right]\nonumber \\
 = &  D_{jlm} + o_{hp}\left(\left(1 + \frac{\|Y_1\|}{\theta} + \frac{\|Y_1\|^2}{\theta^2}\right)\frac{N}{\theta^3} + \frac{\sqrt{N}}{\theta^2}\right) + O_{hp}\left(\left(1 + \frac{\|Y_1\|}{\theta} \right)\frac{1}{\theta^3}\right) + O_{hp}\left(\frac{\sqrt{N}(\log N)^{\xi/2}}{\theta^3}\right) \nonumber\\
&= D_{jlm} + o_{hp}\left( \frac{\sqrt{N}}{\theta^2}\right).\label{EigenVector.3.exten} 
\end{align}
\end{remark}

\begin{lemma}\label{EVec.pivot1}
Let \(1 \leq i \leq k\), and let \(v_i\) denote the normalized eigenvector corresponding to the eigenvalue \(\lambda_i(A)\). Then, as \(N \to \infty\), we have the expansion
\begin{equation}\label{finding}
(e_i^\prime v_i)^{-2} 
= \theta^2\alpha_i^2\, e_i^\prime (\mu I - W)^{-2}e_i \;+\; \Gamma_i \;+\; \Upsilon_i \;+\; \mathcal{E}_N,
\end{equation}
where
\begin{align*}
\Gamma_i &:= 2\theta^3 \mu_0^{-3} \sum_{\substack{1 \le l \le k \\ l \ne i}} (\alpha_l)^{1/2}(\alpha_i)^{3/2}\, \left[Y^{-1} d_{li}\right](l), \\
\Upsilon_i &:= \theta^4 \mu_0^{-4} \sum_{\substack{1 \le l,m \le k \\ l \ne i,\, m \ne i}} (\alpha_l \alpha_m)^{1/2} \alpha_i\, \left[Y^{-1} \tilde d_{lm}\right](l)\, \left[Y^{-1} d_{lm}\right](m), \\
\text{ and, } 
\mathcal{E}_N &:= o_{hp}\left(\frac{\sqrt{N}}{\theta^2}\right)
\end{align*}
denotes an error term which is negligible with high probability.

Here, \(Y\) is a deterministic matrix converging to 
\[
\mathrm{Diag}\left(1-\frac{\alpha_1}{\alpha_i},\ldots,1-\frac{\alpha_{i-1}}{\alpha_i},1-\frac{\alpha_{i+1}}{\alpha_i},\ldots,1-\frac{\alpha_k}{\alpha_i}\right),
\]
and the vectors \(d_{jl}, \tilde d_{jl}\) are deterministic with each component of order \(O\left(\frac{N}{\theta^2}\right)\).
\end{lemma}

\begin{proof}
Let us fix $N\ge1$ and a sample point for which $\|W\|<\mu$, which is a high probability event, as affirmed by Remark~\ref{Bound.W/mu}. The proof proceeds in the following steps.

\textbf{Step 1:}
We begin with the eigenvector equation:
\begin{equation}\label{proof.l1.eq2}
\mu v_i = A v_i = W v_i + \theta \sum_{l=1}^k \alpha_l (e_l^\prime v_i) e_l.
\end{equation}
Since $\|W\|<\mu$, the matrix $\mu I - W$ is invertible, giving
\begin{equation}\label{eq.master1}
v_i = \theta \sum_{l=1}^k \alpha_l (e_l^\prime v_i)(\mu I - W)^{-1} e_l.
\end{equation}

\textbf{Step 2:}
Define the vector
\begin{equation}\label{V.EVec}
u := [\sqrt{\alpha_1}(e_1^\prime v_i), \ldots, \sqrt{\alpha_k}(e_k^\prime v_i)]^\prime.
\end{equation}
Premultiplying \eqref{eq.master1} with $\sqrt{\alpha_j} \mu e_j^\prime$ for $1 \le j \le k$, we find that
\begin{equation}\label{proof.l1.eqnew2}
V u = \mu u,
\end{equation}
where $V(j,l) := \theta \sqrt{\alpha_j \alpha_l} e_j^\prime (\mu I - W)^{-1} e_l$.

\textbf{Step 3:}
Using \eqref{eq.master1}, and normalization of \(v_i\), we expand
\begin{align}
(e_i^\prime v_i)^{-2} &= \theta^2 \alpha_i^2\, e_i^\prime (\mu I - W)^{-2} e_i \label{eq.master2} \\
&\quad + \theta^2 \sum_{\substack{l=i \text{ or } m=i \\ (l,m)\ne(i,i)}} \alpha_l \alpha_m \frac{(e_l^\prime v_i)(e_m^\prime v_i)}{(e_i^\prime v_i)^2} e_l^\prime (\mu I - W)^{-2} e_m \nonumber \\
&\quad + \theta^2 \sum_{\substack{l,m \ne i}} \alpha_l \alpha_m \frac{(e_l^\prime v_i)(e_m^\prime v_i)}{(e_i^\prime v_i)^2} e_l^\prime (\mu I - W)^{-2} e_m. \nonumber
\end{align}

\textbf{Step 4:}
Let $\tilde u$ be the column vector obtained by removing the $i$-th entry from the vector $u$ defined in \eqref{V.EVec}. Let $\tilde V_i$ denote the $i$-th column of $V$ with the $i$-th entry removed, and let $\tilde V$ be the $(k-1) \times (k-1)$ matrix obtained by deleting the $i$-th row and $i$-th column from $V$. Then, equation \eqref{proof.l1.eqnew2} implies that
\begin{equation}\label{Vtruncated1}
\mu \tilde u = \tilde V \tilde u + u(i)\tilde V_i, \qquad \text{w.h.p.}
\end{equation}

By Lemma \ref{EigenVector.1}
\[
\left\| I_{k-1} - \mu^{-1}\tilde V - D \right\| = o_{hp}\left( \frac{\sqrt{N}}{\theta} \right),
\]
where $D = \mathrm{Diag}\left(1 - \frac{\alpha_1}{\alpha_i}, \ldots, 1 - \frac{\alpha_{i-1}}{\alpha_i}, 1 - \frac{\alpha_{i+1}}{\alpha_i}, \ldots, 1 - \frac{\alpha_k}{\alpha_i} \right)$.

Since $D$ is invertible under our ordering assumption $\alpha_1 > \alpha_2 > \cdots > \alpha_k > 0$, it follows that $I_{k-1} - \mu^{-1}\tilde V$ is invertible with high probability. Thus, we may rewrite \eqref{Vtruncated1} as
\begin{equation}\label{Vtruncated2}
\tilde u = u(i)\mu^{-1}\left(I_{k-1} - \mu^{-1} \tilde V\right)^{-1}\tilde V_i, \qquad \text{w.h.p.}
\end{equation}

Substituting \eqref{Vtruncated2} into \eqref{V.EVec}, we find that for any $j \ne i$,
\begin{equation}\label{EigenSp.dim1}
e_j^\prime v_i = \frac{\sqrt{\alpha_i}}{\sqrt{\alpha_j}} e_i^\prime v_i \mu^{-1} \left[ \left(I_{k-1} - \mu^{-1} \tilde V \right)^{-1} \tilde V_i \right](j).
\end{equation}

\textbf{Step 5:}

Define $E := I_{k-1} - \mu^{-1}\tilde V$. The entries of $E$ can be explicitly written in block form as:
\begin{align*}
E_{lm} = 
\begin{cases}
1 - \frac{\theta}{\mu} \alpha_m\, e_m^\prime\left(I - \frac{W}{\mu}\right)^{-1}e_m & \text{if } l = m, \\
- \frac{\theta}{\mu} \sqrt{\alpha_l \alpha_m}\, e_l^\prime\left(I - \frac{W}{\mu}\right)^{-1}e_m & \text{if } l \ne m,
\end{cases}
\end{align*}
with appropriate index shifts for $l,m \ge i$ due to the deletion of the $i$-th row and column.

Each entry $E_{lm}$ can be approximated as:
\begin{align}
E_{lm} 
&= \delta_{lm} - \frac{\theta}{\mu_0} \sqrt{\alpha_l \alpha_m}\, e_l^\prime (\mu I - W)^{-1} e_m \notag \\
&= \delta_{lm} - \frac{\theta}{\mu_0} \sqrt{\alpha_l \alpha_m}(S_2^{lm} + e_l^\prime e_m) + Z_{lm}, \label{Det.Approx1}
\end{align}
where
\[
Z_{lm} = O_{hp}\left( \frac{N|\mu - \mu_0|}{\theta^3} + \frac{\sqrt{N}}{\theta^2} + \frac{\sqrt{N}(\log N)^{\xi/2}}{\theta^2} + \frac{|\mu - \mu_0|}{\theta} + \frac{(\log N)^{\xi/4}}{\theta} \right).
\]

Let $Y$ denote the deterministic matrix with entries
\begin{equation}\label{eq:Ymatrix}
Y_{lm} := \delta_{lm} - \frac{\theta}{\mu_0} \sqrt{\alpha_l \alpha_m} (S_2^{lm} + e_l^\prime e_m).
\end{equation}
Then $E = Y + Z$, and moreover,
\begin{equation}\label{Det.Approx2}
\|Y - D\| = O\left( \frac{N}{\theta^2} \right).
\end{equation}

Let $B_\delta$ be the closed ball (in operator norm) of radius $\delta > 0$ centered at $D$, such that every matrix in $B_\delta$ is invertible. Define the constants
\begin{equation}\label{InverseBound}
C_1 := \sup_{F \in B_\delta} \|F^{-1}\| < \infty, \quad
C_2 := \sup_{\substack{F_1, F_2 \in B_\delta}} \frac{\|F_1^{-1} - F_2^{-1}\|}{\|F_1 - F_2\|} < \infty.
\end{equation}
Then for large $N$, both $Y$ and $I_{k-1} - \mu^{-1}\tilde V$ lie in $B_\delta$ with high probability, and
\begin{equation}
\| (Y + Z)^{-1} - Y^{-1} \| \le C_2 \|Z\|.
\end{equation}
This establishes the deterministic control required to analyze the subsequent terms in equation \eqref{eq.master2}.




\textbf{Step 7:}

Since this analysis occurs on a high-probability event, we define \( \bar V_i := \tilde V_i/\theta \). Using Lemma \ref{EigenVector.1} and the expansion in \eqref{Det.Approx1}, we obtain for \( j \ne i \),
\begin{align}
& \left\| \left[ (I_{k-1} - \mu^{-1} \tilde V)^{-1} \bar V_i - Y^{-1} \bar V_i \right] \cdot e_j^\prime \left(I - \frac{W}{\mu}\right)^{-2} e_i \right\| \notag \\
& \le \left\| (Y + Z)^{-1} - Y^{-1} \right\| \cdot \|\bar V_i\| \cdot \left| e_j^\prime \left(I - \frac{W}{\mu} \right)^{-2} e_i \right| \notag \\
& = o_{hp} \left( \left(1 + \frac{\|Y_1\|}{\theta} \right) \frac{N}{\theta^3} + \frac{\sqrt{N}}{\theta^2} \right)=:\widetilde{\mathcal{E}_N}, \label{EVec.KeyApprox0}
\end{align}
where the final bound follows from \eqref{eq:mu-mu0}.

Thus, for \( j \ne i \), using \eqref{EVec.KeyApprox0} and \eqref{EigenSp.dim1} we deduce for the first term in the series \eqref{eq.master2}:
\begin{align}
\theta^2 \frac{e_j^\prime v_i}{e_i^\prime v_i} \cdot e_j^\prime (\mu I - W)^{-2} e_i 
= \frac{\sqrt{\alpha_i}}{\sqrt{\alpha_j}}\theta^3 \mu^{-3} (Y^{-1} \bar V_i)(j) \cdot e_j^\prime \left(I - \frac{W}{\mu}\right)^{-2} e_i + \widetilde{\mathcal{E}_N}. \label{EVec.KeyApprox1}
\end{align}

\textbf{Step 8:}

Applying the above analysis again for \( l,m \ne i \), we get:
\begin{align}
& \theta^2 \frac{e_l^\prime v_i}{e_i^\prime v_i} \cdot \frac{e_m^\prime v_i}{e_i^\prime v_i} \cdot e_l^\prime (\mu I - W)^{-2} e_m \notag \\
= & \frac{\alpha_i}{\sqrt{\alpha_l}\sqrt{\alpha_m}}\theta^4 \mu^{-4} (Y^{-1} \bar V_i)(l) (Y^{-1} \bar V_i)(m) \cdot e_l^\prime \left(I - \frac{W}{\mu}\right)^{-2} e_m + \widetilde{\mathcal{E}_N}. \label{EVec.KeyApprox2}
\end{align}

\textbf{Step 9:}

Let \( d_{ji} \in \mathbb{R}^{k-1} \) denote the deterministic vector defined via the entries \( D_{mji} \) from \eqref{EigenVector.3.intermd}. Using \eqref{Det.Approx2} and \eqref{EigenVector.3.exten}, we have \( \|Y^{-1}\| = O(1) \), and Lemma \ref{EigenVector.3} gives
\begin{align}
\left\| Y^{-1} \left( \bar V_i e_j^\prime (\mu I - W)^{-2} e_i - d_{ji} \right) \right\| 
= o_{hp} \left( \left(1 + \frac{\|Y_1\|}{\theta} + \frac{\|Y_1\|^2}{\theta^2} \right) \frac{N}{\theta^3} + \frac{\sqrt{N}}{\theta^2} \right).
\end{align}

Thus, \eqref{EVec.KeyApprox1} becomes
\begin{align}
\theta^2 \frac{e_j^\prime v_i}{e_i^\prime v_i} \cdot e_j^\prime (\mu I - W)^{-2} e_i 
= \frac{\sqrt{\alpha_i}}{\sqrt{\alpha_j}}\theta^3 \mu_0^{-3} (Y^{-1} d_{ji})(j) + \mathcal{E}_N. \label{EVec.KeyApprox3}
\end{align}

\textbf{Step 10:}

Define \( \tilde d_{lm} \in \mathbb{R}^{k-1} \) with $(\tilde d_{lm})_{j}= D_{jlm} S_2^{j^\prime i}$ for $j\neq i$ as in Remark \ref{Remark:djlm}. Arguing as before, we find
\begin{align}
\left\| Y^{-1}(\bar V_i - \tilde d_{lm}) Y^{-1} d_{lm} \right\| = \mathcal{E}_N.
\end{align}
Hence, from \eqref{EVec.KeyApprox2} we conclude:
\begin{align}
\theta^2 \frac{e_l^\prime v_i}{e_i^\prime v_i} \cdot \frac{e_m^\prime v_i}{e_i^\prime v_i} \cdot e_l^\prime (\mu I - W)^{-2} e_m 
=\frac{\alpha_i}{\sqrt{\alpha_l}\sqrt{\alpha_m}} \theta^4 \mu_0^{-4} (Y^{-1} \tilde d_{lm})(l) (Y^{-1} d_{lm})(m) + \mathcal{E}_N. \label{EVec.KeyApprox4}
\end{align}


Hence substituting the approximations \eqref{EVec.KeyApprox3} and \eqref{EVec.KeyApprox4} into equation \eqref{eq.master2}, we conclude the desired expansion:
\[
(e_i^\prime v_i)^{-2} = \theta^2 \alpha_i^2 e_i^\prime (\mu I - W)^{-2} e_i + \Gamma_i + \Upsilon_i + \mathcal{E}_N,
\]
as stated in Lemma \ref{EVec.pivot1}.
\end{proof}

\end{document}
